% Options for packages loaded elsewhere
% Options for packages loaded elsewhere
\PassOptionsToPackage{unicode}{hyperref}
\PassOptionsToPackage{hyphens}{url}
\PassOptionsToPackage{dvipsnames,svgnames,x11names}{xcolor}
%
\documentclass[
  a4paper,
  DIV=11,
  numbers=noendperiod]{scrartcl}
\usepackage{xcolor}
\usepackage[margin=1in]{geometry}
\usepackage{amsmath,amssymb}
\setcounter{secnumdepth}{5}
\usepackage{iftex}
\ifPDFTeX
  \usepackage[T1]{fontenc}
  \usepackage[utf8]{inputenc}
  \usepackage{textcomp} % provide euro and other symbols
\else % if luatex or xetex
  \usepackage{unicode-math} % this also loads fontspec
  \defaultfontfeatures{Scale=MatchLowercase}
  \defaultfontfeatures[\rmfamily]{Ligatures=TeX,Scale=1}
\fi
\usepackage{lmodern}
\ifPDFTeX\else
  % xetex/luatex font selection
\fi
% Use upquote if available, for straight quotes in verbatim environments
\IfFileExists{upquote.sty}{\usepackage{upquote}}{}
\IfFileExists{microtype.sty}{% use microtype if available
  \usepackage[]{microtype}
  \UseMicrotypeSet[protrusion]{basicmath} % disable protrusion for tt fonts
}{}
\makeatletter
\@ifundefined{KOMAClassName}{% if non-KOMA class
  \IfFileExists{parskip.sty}{%
    \usepackage{parskip}
  }{% else
    \setlength{\parindent}{0pt}
    \setlength{\parskip}{6pt plus 2pt minus 1pt}}
}{% if KOMA class
  \KOMAoptions{parskip=half}}
\makeatother
% Make \paragraph and \subparagraph free-standing
\makeatletter
\ifx\paragraph\undefined\else
  \let\oldparagraph\paragraph
  \renewcommand{\paragraph}{
    \@ifstar
      \xxxParagraphStar
      \xxxParagraphNoStar
  }
  \newcommand{\xxxParagraphStar}[1]{\oldparagraph*{#1}\mbox{}}
  \newcommand{\xxxParagraphNoStar}[1]{\oldparagraph{#1}\mbox{}}
\fi
\ifx\subparagraph\undefined\else
  \let\oldsubparagraph\subparagraph
  \renewcommand{\subparagraph}{
    \@ifstar
      \xxxSubParagraphStar
      \xxxSubParagraphNoStar
  }
  \newcommand{\xxxSubParagraphStar}[1]{\oldsubparagraph*{#1}\mbox{}}
  \newcommand{\xxxSubParagraphNoStar}[1]{\oldsubparagraph{#1}\mbox{}}
\fi
\makeatother


\usepackage{longtable,booktabs,array}
\newcounter{none} % for unnumbered tables
\usepackage{calc} % for calculating minipage widths
% Correct order of tables after \paragraph or \subparagraph
\usepackage{etoolbox}
\makeatletter
\patchcmd\longtable{\par}{\if@noskipsec\mbox{}\fi\par}{}{}
\makeatother
% Allow footnotes in longtable head/foot
\IfFileExists{footnotehyper.sty}{\usepackage{footnotehyper}}{\usepackage{footnote}}
\makesavenoteenv{longtable}
\usepackage{graphicx}
\makeatletter
\newsavebox\pandoc@box
\newcommand*\pandocbounded[1]{% scales image to fit in text height/width
  \sbox\pandoc@box{#1}%
  \Gscale@div\@tempa{\textheight}{\dimexpr\ht\pandoc@box+\dp\pandoc@box\relax}%
  \Gscale@div\@tempb{\linewidth}{\wd\pandoc@box}%
  \ifdim\@tempb\p@<\@tempa\p@\let\@tempa\@tempb\fi% select the smaller of both
  \ifdim\@tempa\p@<\p@\scalebox{\@tempa}{\usebox\pandoc@box}%
  \else\usebox{\pandoc@box}%
  \fi%
}
% Set default figure placement to htbp
\def\fps@figure{htbp}
\makeatother





\setlength{\emergencystretch}{3em} % prevent overfull lines

\providecommand{\tightlist}{%
  \setlength{\itemsep}{0pt}\setlength{\parskip}{0pt}}



 
\usepackage[]{biblatex}
\addbibresource{references.bib}


\KOMAoption{captions}{tableheading}
\makeatletter
\@ifpackageloaded{caption}{}{\usepackage{caption}}
\AtBeginDocument{%
\ifdefined\contentsname
  \renewcommand*\contentsname{Table of contents}
\else
  \newcommand\contentsname{Table of contents}
\fi
\ifdefined\listfigurename
  \renewcommand*\listfigurename{List of Figures}
\else
  \newcommand\listfigurename{List of Figures}
\fi
\ifdefined\listtablename
  \renewcommand*\listtablename{List of Tables}
\else
  \newcommand\listtablename{List of Tables}
\fi
\ifdefined\figurename
  \renewcommand*\figurename{Figure}
\else
  \newcommand\figurename{Figure}
\fi
\ifdefined\tablename
  \renewcommand*\tablename{Table}
\else
  \newcommand\tablename{Table}
\fi
}
\@ifpackageloaded{float}{}{\usepackage{float}}
\floatstyle{ruled}
\@ifundefined{c@chapter}{\newfloat{codelisting}{h}{lop}}{\newfloat{codelisting}{h}{lop}[chapter]}
\floatname{codelisting}{Listing}
\newcommand*\listoflistings{\listof{codelisting}{List of Listings}}
\makeatother
\makeatletter
\makeatother
\makeatletter
\@ifpackageloaded{caption}{}{\usepackage{caption}}
\@ifpackageloaded{subcaption}{}{\usepackage{subcaption}}
\makeatother
\usepackage{bookmark}
\IfFileExists{xurl.sty}{\usepackage{xurl}}{} % add URL line breaks if available
\urlstyle{same}
\hypersetup{
  pdftitle={Architecting a Business with MLOps and Data Science},
  pdfauthor={N. Javier Buitrago Aza},
  pdfkeywords={Business Architecture, MLOps, Data Science, Enterprise
Learning Systems, Quality Attributes, Experimentation},
  colorlinks=true,
  linkcolor={blue},
  filecolor={Maroon},
  citecolor={Blue},
  urlcolor={Blue},
  pdfcreator={LaTeX via pandoc}}


\title{Architecting a Business with MLOps and Data Science}
\usepackage{etoolbox}
\makeatletter
\providecommand{\subtitle}[1]{% add subtitle to \maketitle
  \apptocmd{\@title}{\par {\large #1 \par}}{}{}
}
\makeatother
\subtitle{From Models to Enterprise Learning Systems}
\author{N. Javier Buitrago Aza}
\date{2026-06-04}
\begin{document}
\maketitle
\begin{abstract}
Most organizations invest heavily in Data Science, Machine Learning and
MLOps, yet many struggle to turn analytical insights into measurable
business value. This article argues that Data Science and MLOps should
be treated as components of a broader enterprise architecture and that
modern organizations must evolve from model-centric thinking toward
\emph{enterprise learning systems}: architectures capable of
continuously sensing, learning, adapting and generating business value.
\end{abstract}

\renewcommand*\contentsname{Table of contents}
{
\setcounter{tocdepth}{3}
\tableofcontents
}

\section{Introduction}\label{introduction}

Most organizations invest heavily in Data Science, Machine Learning and
MLOps across the data engineering, data governance and platform layers.
In the particular case of models --- classical classifiers, regressions,
even forecasting models --- models are trained, dashboards are built,
experiments are executed, and cloud platforms process vast amounts of
data every day. Yet despite these investments, many organizations still
struggle to convert analytical insights into measurable business value.

One can ask what happens with the expected results, and the answer is
usually the same: \textbf{the problem is not technological, the problem
is architectural.}

To be concrete, machine learning models are often treated as isolated
technical assets rather than as components of a larger enterprise
system. Accuracy improves, yet decision quality remains unchanged.
Experiments produce statistically significant results, yet business
outcomes fail to improve. Monitoring platforms detect drift, yet
strategic decisions continue to rely on outdated assumptions.

This article argues that Data Science and MLOps should be viewed as
components of a broader enterprise architecture. Modern organizations
must evolve from model-centric thinking toward \textbf{enterprise
learning systems}: architectures capable of continuously sensing,
learning, adapting and generating business value. Generative AI is, of
course, also a standard to consider; for the purposes of this paper the
same discussions apply.

\section{The Problem with Traditional Data
Science}\label{the-problem-with-traditional-data-science}

Most machine learning initiatives optimize \emph{models} instead of
optimizing the \emph{enterprise}. Teams frequently focus on:

\begin{itemize}
\tightlist
\item
  Prediction accuracy.
\item
  AUC and F1 scores.
\item
  \(R^2\), RMSE and similar regression metrics.
\item
  Feature engineering.
\item
  Hyperparameter tuning.
\end{itemize}

These metrics are important, but they are not business outcomes.

A highly accurate model with low adoption, poor operational integration
or limited strategic impact generates little value. Business value
emerges only when \textbf{predictions become decisions},
\textbf{decisions become actions} and \textbf{actions influence
measurable outcomes}.

This observation leads to a simple but powerful principle:

\begin{quote}
The objective of Machine Learning is not prediction. The objective is
\textbf{improved decision-making}.
\end{quote}

\section{The Enterprise as a (Dynamic) Learning
System}\label{the-enterprise-as-a-dynamic-learning-system}

A modern enterprise should not be understood as a static collection of
processes, applications, dashboards and machine learning models. It is
better understood as a \textbf{dynamic system}: its state changes over
time as leaders launch experiments, deploy models, modify processes,
reallocate resources and respond to market signals.

It is also a \textbf{learning system}, because those actions generate
outcomes and those outcomes should feed back into future decisions. In
this sense, MLOps is not only a technical discipline for deploying
models; it is part of the enterprise feedback mechanism that allows the
organization to observe, act, measure and adapt.

\subsection{A formal model of the
enterprise}\label{a-formal-model-of-the-enterprise}

An enterprise \(\mathcal{E}\) can be viewed as a dynamic system composed
of capabilities \(\mathcal{C}\), processes \(\mathcal{P}\), resources
\(\mathcal{R}\), data assets \(\mathcal{D}\) and decision mechanisms
\(\mathcal{M}\). In the context of MLOps and Data Science, the decision
mechanisms of interest are machine learning models; for the purposes of
this discussion we therefore denote the set of all ML models by
\(\mathcal{M}\).

Let \(\mathcal{S}\) be the set of possible enterprise states, let
\(\mathcal{A}\) be the set of possible organizational actions, and let
\(t \in \mathbb{N} \cup \{0\}\) denote discrete time. A
\textbf{state-transition function} is a map

\[
F \colon \mathcal{S} \times \mathcal{A} \rightarrow \mathcal{S},
\]

such that

\[
S_{t+1} = F(S_t, A_t).
\]

Here \(S_t \in \mathcal{S}\) represents the enterprise state at time
\(t\), and \(A_t \in \mathcal{A}\) represents the organizational action
taken at time \(t\). The state \(S_t\) may include the enterprise
configuration at that time,

\[
S_t \;\equiv\; (C_t, P_t, R_t, D_t, M_t),
\]

where

\[
C_t \subset \mathcal{C}, \quad
P_t \subset \mathcal{P}, \quad
R_t \subset \mathcal{R}, \quad
D_t \subset \mathcal{D}, \quad
M_t \subset \mathcal{M}.
\]

The initial state \(S_0\) represents the initial enterprise
configuration,

\[
S_0 \;\equiv\; (C_0, P_0, R_0, D_0, M_0).
\]

The transition equation \(S_{t+1} = F(S_t, A_t)\) therefore expresses
that the next enterprise state is determined by the current enterprise
state and the action applied to it.

\subsection{The enterprise value
function}\label{the-enterprise-value-function}

We can now define an \textbf{enterprise value function} as a map

\[
V \colon \mathcal{S} \rightarrow \mathbb{R},
\]

where \(V(S_t)\) represents the business value associated with the
enterprise state \(S_t \in \mathcal{S}\) at time
\(t \in \mathbb{N} \cup \{0\}\). The sequence
\(\{V(S_t)\}_{t \in \mathbb{N} \cup \{0\}}\) therefore represents the
realized enterprise value over time.

The value generated by the action \(A_t\) is defined as

\[
\Delta V_t \coloneqq V(S_{t+1}) - V(S_t),
\]

and using the transition function,

\[
\Delta V_t = V(F(S_t, A_t)) - V(S_t).
\]

From a business perspective, an action \(A_t\) is:

\begin{itemize}
\tightlist
\item
  \textbf{value-increasing} if \(\Delta V_t > 0\),
\item
  \textbf{value-neutral} if \(\Delta V_t = 0\),
\item
  \textbf{value-destroying} if \(\Delta V_t < 0\).
\end{itemize}

In the MLOps and Data Science context, \(A_t\) may represent deploying a
model, launching an experiment, retraining a model, changing a campaign,
reallocating budget or modifying an operational process. The formulation
is intentionally general: it does not assume that a model improves the
enterprise by default. \textbf{A model creates value only if the
transition it induces moves the enterprise from \(S_t\) to \(S_{t+1}\)
in a way that increases \(V\).}

\section{Business Architecture as the Missing
Layer}\label{business-architecture-as-the-missing-layer}

Business Architecture provides the structure that connects strategy,
operations and technology. Without it:

\begin{itemize}
\tightlist
\item
  Models remain disconnected from strategy.
\item
  Experiments remain disconnected from operations.
\item
  Metrics remain disconnected from governance.
\end{itemize}

Business \textbf{capabilities} become the primary mechanism through
which analytical insights create enterprise value. Instead of asking

\begin{quote}
\emph{What model should we build?}
\end{quote}

organizations should ask

\begin{quote}
\emph{What business capability are we trying to improve?}
\end{quote}

This shift fundamentally changes how MLOps systems are designed and
governed.

\section{Architecturally Significant
Requirements}\label{architecturally-significant-requirements}

In software architecture, \textbf{Architecturally Significant
Requirements (ASRs)} are requirements that strongly influence
architectural decisions. For MLOps systems, ASRs typically focus on
\textbf{quality attributes} rather than predictive outputs.
Representative examples include:

\begin{itemize}
\tightlist
\item
  \textbf{Scalability.} Real-time, consistent predictions over a given
  time interval without degradation of the system as the request volume
  grows.
\item
  \textbf{Availability.} The system must remain operational under faults
  and load spikes, with explicit recovery-time and recovery-point
  objectives.
\item
  \textbf{Security.} End-to-end protection of data and models, including
  authentication, authorization, encryption in transit and at rest, and
  intrusion-detection coverage of the inference path.
\item
  \textbf{Observability.} The ability to inspect, trace and explain
  every prediction, every retraining and every data transformation in
  production.
\item
  \textbf{Modifiability.} Changes to features, models or pipelines can
  be introduced and rolled back safely without disrupting downstream
  consumers.
\item
  \textbf{Explainability.} Predictions can be explained to business
  stakeholders and regulators in terms aligned with the decisions they
  support.
\end{itemize}

A useful mechanism for discovering ASRs is the \textbf{Utility Tree},
which decomposes stakeholder concerns into measurable quality scenarios.
Architectures can then be evaluated using methods such as the
\textbf{Architecture Tradeoff Analysis Method (ATAM)}, exposing risks,
sensitivities and architectural tradeoffs before systems reach
production.

\section{Experimentation as a Business
Capability}\label{experimentation-as-a-business-capability}

Experimentation is not merely a statistical activity. It is a
\textbf{strategic business capability}. The fundamental causal quantity
of interest is the \textbf{Average Treatment Effect},

\[
ATE = \mathbb{E}[Y(1) - Y(0)],
\]

where \(Y(1)\) represents the outcome under treatment and \(Y(0)\)
represents the outcome under control. In practice, organizations
estimate

\[
\widehat{ATE} = \bar{Y}_T - \bar{Y}_C.
\]

Before evaluating business impact, however, \textbf{experiment quality
must be validated}. A common balance diagnostic is the
\textbf{Standardized Mean Difference},

\[
SMD(X) = \frac{\bar{X}_T - \bar{X}_C}{\sqrt{\dfrac{s_T^2 + s_C^2}{2}}}.
\]

Typically \(|SMD| < 0.1\) indicates acceptable covariate balance.

For monetary outcomes, which frequently exhibit unequal variances,
\textbf{Welch's t-test} provides a more robust inference framework:

\[
t_W = \frac{\bar{Y}_T - \bar{Y}_C}{\sqrt{\dfrac{s_{Y,T}^2}{n_T}
                                          + \dfrac{s_{Y,C}^2}{n_C}}}.
\]

A mature experimentation capability therefore evaluates both
\textbf{experimental validity} and \textbf{business impact}.

\section{Propensity Models as Strategic Decision
Engines}\label{propensity-models-as-strategic-decision-engines}

Propensity models estimate the probability that a customer or entity
performs a target action. Formally,

\[
p_i = \mathbb{P}(Y_i = 1 \mid X_i).
\]

Examples include:

\begin{itemize}
\tightlist
\item
  Purchase propensity.
\item
  Churn propensity.
\item
  Upgrade propensity.
\item
  Response propensity.
\end{itemize}

These models become strategic instruments because they allocate scarce
resources toward the most valuable opportunities. A common operational
metric is the \textbf{Decile Lift},

\[
Lift_{D_1} = \frac{\mathbb{E}[Y \mid D_1]}{\mathbb{E}[Y]}.
\]

High lift indicates strong prioritization capability and efficient
resource allocation.

\section{Drift as Enterprise Risk}\label{drift-as-enterprise-risk}

Models degrade over time because businesses evolve. Two common
mechanisms are:

\begin{itemize}
\tightlist
\item
  \textbf{Covariate shift.}
\item
  \textbf{Concept drift.}
\end{itemize}

To monitor changing populations, organizations frequently employ the
\textbf{Population Stability Index},

\[
PSI = \sum_{k=1}^{K} (A_k - E_k)\,
      \ln\!\left( \frac{A_k}{E_k} \right).
\]

A common interpretation is:

{\def\LTcaptype{none} % do not increment counter
\begin{longtable}[]{@{}ll@{}}
\toprule\noalign{}
\(PSI\) value & Interpretation \\
\midrule\noalign{}
\endhead
\bottomrule\noalign{}
\endlastfoot
\(PSI < 0.10\) & Stable \\
\(0.10 \le PSI < 0.25\) & Moderate shift \\
\(PSI \ge 0.25\) & Significant shift \\
\end{longtable}
}

Drift should not be viewed solely as a technical concern; it is an
\textbf{enterprise risk indicator}. We may define the \textbf{Enterprise
Machine Learning Risk} as

\[
EMLR = P_{\text{drift}} \times \text{Impact}_{\text{business}}.
\]

This transforms monitoring from a technical activity into a governance
mechanism.

\section{Stratified Modeling at
Scale}\label{stratified-modeling-at-scale}

Many organizations operate across thousands of customers, stores,
products, routes, vessels or geographic regions. A single global model
may fail to capture these differences. Instead, organizations deploy
\textbf{model families},

\[
\mathcal{M} \coloneqq \{m_1, m_2, \ldots, m_N\},
\]

so that predictions become

\[
\hat{y}_i = m_i(x_i).
\]

This architectural pattern is known as \textbf{stratified modeling}.
Each model requires:

\begin{itemize}
\tightlist
\item
  Independent monitoring.
\item
  Dedicated thresholds.
\item
  Retraining policies.
\item
  Governance controls.
\end{itemize}

Consequently, scalable MLOps architectures must support massively
parallel training and deployment pipelines.

\section{The Strategic Recommendation
Layer}\label{the-strategic-recommendation-layer}

Data alone does not create business value: insights must be transformed
into decisions. The \textbf{Strategic Recommendation Layer} converts
analytical evidence into governance actions. A typical decision sequence
evaluates:

\begin{enumerate}
\def\labelenumi{\arabic{enumi}.}
\tightlist
\item
  Statistical power.
\item
  Experimental validity.
\item
  Effect size.
\item
  Model trustworthiness.
\item
  Operational feasibility.
\end{enumerate}

The resulting diagnoses map directly to recommendations:

{\def\LTcaptype{none} % do not increment counter
\begin{longtable}[]{@{}ll@{}}
\toprule\noalign{}
Diagnosis & Recommendation \\
\midrule\noalign{}
\endhead
\bottomrule\noalign{}
\endlastfoot
Significant positive impact & Scale \\
High drift & Retrain \\
Low power & Collect additional data \\
Low lift & Recalibrate targeting \\
Negative impact & Retire model \\
\end{longtable}
}

Analytics therefore becomes a \textbf{governance capability} rather than
a reporting capability.

\section{Conclusion}\label{conclusion}

The future of Data Science is not the model. The future is the
\textbf{enterprise learning system}.

Organizations that successfully integrate Business Architecture,
experimentation, MLOps, governance and continuous feedback mechanisms
create adaptive systems capable of evolving alongside the business
itself. In such organizations, machine learning is no longer an isolated
technical discipline: it becomes a \textbf{strategic capability embedded
within the architecture of the enterprise}.


\printbibliography



\end{document}
